% SPDX-License-Identifier: CC-BY-4.0
% Copyright 2025 Florian Aram Feuerriegel — kassensturz.org
% Teil III — Das Modell als Wirkungsgraph
%
% Alle Zahlen in Kanten, Tabellen und Kurven stammen aus generiert/ (messung.js).

\section{Wie die Graphen entstehen}
\label{sec:graphmethode}

Teil~II beschreibt das Modell als Gleichungssystem. Man kann dasselbe System als
\emph{gerichteten, gewichteten Graphen} lesen: Knoten sind Größen, eine Kante $u\to v$ heißt
\enquote{$u$ geht in die Berechnung von $v$ ein}, und ihr Gewicht sagt, um wie viel sich $v$
ändert, wenn $u$ sich um eine Einheit ändert. Der Graph hat drei Ebenen:

\begin{enumerate}
\item \textbf{Stellgrößen} (die Regler am Verhandlungstisch),
\item \textbf{Kanäle} (Zwischengrößen: Aufkommen je Steuerart, Transfers, Renten,
  Arbeitsangebot, Investitionen, Nachfrage),
\item \textbf{Kennzahlen} (was das Spiel bewertet: Saldo, Verteilung, Emissionen, BIP,
  Schulden).
\end{enumerate}

Die Gewichte stammen aus zwei Quellen, und der Unterschied ist wichtig:

\begin{itemize}
\item \textbf{Stellgröße $\to$ Kanal: gemessen.} Jede Stellgröße wurde, ausgehend vom Status
  quo, um einen festen Schritt verschoben (Tabelle~\ref{tab:schritte}), die Engine rechnete
  die erste Periode (2025--2028) neu, und die Differenz jedes Kanals zum Status quo ist das
  Gewicht. Formal ist das die endliche Differenz
  \begin{equation}
    w_{u\to v}=\frac{v(\theta^{SQ}+h\,e_u)-v(\theta^{SQ})}{1}\qquad\text{je Schritt }h,
  \end{equation}
  keine Ableitung: Die Engine hat Knicke (Abschnitt~\ref{sec:eigenschaften}), und die
  Wirkung gilt genau für den angegebenen Schritt. Wo das Modell nichtlinear ist, zeigen es
  die Kurven in Abschnitt~\ref{sec:kurven}.
\item \textbf{Kanal $\to$ Kennzahl: strukturell.} Diese Gewichte folgen direkt aus den
  Gleichungen von Teil~II, ausgewertet im Status quo. Ein Euro zusätzliches Aufkommen
  verbessert den Saldo um einen Euro; ein Euro Nachfrage bringt \Wstabil\,\euro{} Einnahmen.
\end{itemize}

Die Summe über alle Pfade von einer Stellgröße zu einer Kennzahl ergibt deren
Gesamtwirkung. Wegen der Rückkopplung über die Nachfrage ist der Graph nicht kreisfrei; die
Engine löst den Kreis, indem sie die Nachfrage einmal gegen den Status quo derselben Periode
rechnet (keine Iteration bis zum Fixpunkt).

\begin{table}[!htb]
\centering\small
\begin{tabular}{lll}
\toprule
Art der Stellgröße & Schritt $h$ & Beispiele \\
\midrule
Steuer- und Beitragssätze & $+1$\,\Pp & Eingangs-, Spitzen-, Umsatz-, Körperschaftsteuer, RV, KV \\
Rentenniveau & $+1$\,\Pp & 48\,\% $\to$ 49\,\% \\
Grundfreibetrag & $+1.000$\,\euro & 12.348 $\to$ 13.348\,\euro \\
Monatliche Leistungen & $+10$\,\euro/Monat & Bürgergeld, Kindergeld \\
\COz-Preis & $+10$\,\euro/t & 55 $\to$ 65\,\euro/t \\
Schalter & aus $\to$ ein & Klimageld \\
\bottomrule
\end{tabular}
\caption{Schritte der Messung. Die Stellgrößen der klassischen Fläche stehen mit ihren
Schritten in Tabelle~\ref{tab:zerlegung}.}
\label{tab:schritte}
\end{table}

\section{Der Strukturgraph}
\label{sec:struktur}

Abbildung~\ref{fig:struktur} zeigt, wie die Kanäle auf die Kennzahlen wirken, unabhängig
davon, welche Stellgröße sie bewegt hat. Er ist das Gerüst, in das sich die gemessenen
Ressortgraphen einhängen.

\begin{figure}[p]
\centering
\resizebox{\textwidth}{!}{%
\begin{tikzpicture}[
  kn/.style={draw=grau, rounded corners=2pt, fill=hell, align=center, font=\small, text width=27mm, minimum height=10mm},
  zk/.style={draw=black, thick, rounded corners=2pt, fill=white, align=center, font=\small, text width=27mm, minimum height=10mm},
  pf/.style={-{Stealth[length=2.2mm]}, thick},
  rk/.style={-{Stealth[length=2.2mm]}, thick, dashed, grau},
  gl/.style={font=\scriptsize, fill=white, inner sep=1pt, align=center}]
% Spalte A: Kanäle
\node[kn] (aus)  at (0, 1.9)   {Ausgaben an Haushalte\\{\scriptsize Transfers, Renten}};
\node[kn] (ein)  at (0,-0.3)   {Einnahmen\\{\scriptsize ESt, USt, KSt, GewSt,\\Verm., \COz, Beiträge}};
\node[kn] (arb)  at (0,-5.8)   {Arbeitsangebot $\bar\lambda$};
\node[kn] (inv)  at (0,-7.7)   {Investitionsfaktor $I_f$};
\node[kn] (emi)  at (0,-9.8)   {Emissionen $E$};
% Spalte B: Haushalte und Nachfrage
\node[kn] (hh)   at (4.3,-2.4) {Nettoeinkommen\\der Haushalte $\Delta_i$};
\node[kn] (nf)   at (4.3,-4.7) {Nachfrageimpuls $J$};
% Spalte C: Kennzahlen der Periode und Niveau
\node[zk] (sal)  at (8.6, 0.8) {Saldo $S$};
\node[zk] (ver)  at (8.6,-2.4) {Gini, Palma,\\Armutsrisiko\\{\scriptsize auf Äquivalenzeinkommen}};
\node[zk] (bip)  at (8.6,-4.7) {BIP der Periode\\$Y\eta$};
\node[zk] (bipl) at (8.6,-6.9) {BIP-Niveau\\Folgeperiode};
\node[zk] (cob)  at (8.6,-9.8) {\COz-Budget\\(kumuliert)};
% Spalte D: Schulden
\node[zk] (sq)   at (12.9,-2.4) {Schuldenquote\\Folgeperiode};
% Kanten der Periode
\draw[pf] (ein.east) -- node[gl, pos=0.5]{$+1$} (sal.west);
\draw[pf] (aus.east) -- node[gl, pos=0.75]{$-1$} (sal.west);
\draw[pf] (ein.east) -- node[gl, pos=0.5, below left=0pt and 1pt]{$-1$ {\tiny Inzidenz}} (hh.west);
\draw[pf] (aus.east) -- node[gl, pos=0.18]{$+1$} (hh.west);
\draw[pf] (hh.east)  -- (ver.west);
\draw[pf] (hh.south) -- node[gl]{$\frac{0{,}6}{0{,}48}\,m_i$ {\tiny(0,19 bis 0,81)}} (nf.north);
\draw[pf] (nf.east)  -- node[gl]{$+1$} (bip.west);
\draw[pf] (nf.north east) -- node[gl, pos=0.78]{$+\Wstabil$} (sal.south west);
% Kanten über die Periode hinaus
\draw[pf] (sal.east) -- node[gl]{$\Wschuld$\,\Pp\\{\tiny je Mrd./a}} (sq.north west);
\draw[pf] (bipl.east) -- node[gl, pos=0.5]{Nenner} (sq.south);
\draw[pf] (arb.east) -- node[gl, pos=0.5]{$+\Warbeit$ {\tiny Mrd. je \%}} (bipl.west);
\draw[pf] (inv.east) -- node[gl, pos=0.5]{$+\Winvestn$ / $\Winvestlang$ {\tiny Mrd. je \%, 1 Per. / lang}} (bipl.south west);
\draw[pf] (emi.east) -- node[gl]{$\times\,n$ Jahre} (cob.west);
% Rückwirkungen (gestrichelt)
\draw[rk] (bipl.north) -- node[gl, text=grau]{{\tiny Ausgangsniveau}\\{\tiny nächste Periode}} (bip.south);
\draw[rk] (arb.west) to[out=160,in=200,looseness=1.0] node[gl, pos=0.5, text=grau]{Basis} (ein.west);
\draw[rk] (inv.west) to[out=170,in=190,looseness=1.3] node[gl, pos=0.5, text=grau]{Basis $\pi$} (ein.south west);
\draw[rk] (emi.west) to[out=175,in=185,looseness=1.5] node[gl, pos=0.5, text=grau]{\COz-Basis} ($(ein.south west)+(0,0.2)$);
\end{tikzpicture}}
\caption{Strukturgraph: wie Kanäle auf Kennzahlen wirken. Durchgezogene Kanten tragen das
Gewicht aus den Gleichungen von Teil~II im Status quo: \euro{} Kennzahl je \euro{} Kanal,
wenn nicht anders angegeben. Gestrichelte Kanten sind Rückwirkungen auf die
Bemessungsgrundlagen (Verhalten) und in die nächste Periode. $m_i$ ist die
Grenzkonsumneigung (0,15 in D10c bis 0,65 in D1).}
\label{fig:struktur}
\end{figure}

\paragraph{So liest man den Strukturgraphen.} Beginnen Sie bei einem Kanal und folgen Sie
den Pfeilen. Beispiel \emph{ein Euro mehr Rente}: Er verschlechtert den Saldo um einen Euro
(Kante \enquote{Ausgaben $\to$ Saldo}, $-1$) und erhöht das Nettoeinkommen eines Haushalts um
einen Euro ($+1$). Dieser Haushalt gibt davon je nach Dezil 15 bis 65~Cent aus; mit dem
Multiplikator werden daraus 19 bis 81~Cent Nachfrage. Jeder Euro Nachfrage hebt das BIP der
Periode um einen Euro und bringt \Wstabil\,\euro{} Einnahmen zurück (Kante
\enquote{Nachfrage $\to$ Saldo}). Ein Euro Rente an einen Haushalt aus D1 kostet den Staat
also nicht einen Euro, sondern rund $1-0{,}81\cdot\Wstabil\approx0{,}70$\,\euro. Über die
Periode hinaus wirkt nur, was die gestrichelten und die unteren Kanten tragen: Arbeit,
Kapital und die Schulden.

\section{Die Ressortgraphen}
\label{sec:ressortgraphen}

Die folgenden vier Abbildungen zeigen für jedes Ressort am Verhandlungstisch, welche Kanäle
seine Stellgrößen bewegen. Die Kanäle stehen in allen vier Graphen an derselben Stelle,
damit man sie vergleichen kann.

\paragraph{So liest man einen Ressortgraphen.}
\begin{itemize}
\item \textbf{Links} steht die Stellgröße mit ihrem Schritt und ihrer Gesamtwirkung auf den
  Saldo der ersten Periode, \textbf{rechts} die Kanäle mit ihrer Einheit.
\item Die \textbf{Zahl an der Kante} ist die gemessene Änderung des Kanals, in seiner
  Einheit, für genau diesen Schritt.
\item \textbf{Blau} heißt: Der Kanal steigt. \textbf{Orange} heißt: Er sinkt. Ob das gut ist,
  hängt vom Kanal ab: Steigende Transfers sind für den Saldo schlecht, für die Haushalte gut.
\item Die \textbf{Strichstärke} ist proportional zum Betrag, gemessen am größten Wert
  derselben Einheit über alle Ressorts. Dicke Kanten sind also in allen vier Graphen gleich
  groß.
\item Kanten unter 0,3\,\Mrd{} beziehungsweise 0,02\,\% sind weggelassen.
\end{itemize}
Beispiel aus Abbildung~\ref{fig:soziales}: Ein Punkt Rentenbeitrag bringt
\wert{kanal}{rv}{sv}\,\Mrd{} Beiträge (dicke blaue Kante), senkt aber das Aufkommen aus
Umsatzsteuer um \wert{kanal}{rv}{mwst}\,\Mrd{} und aus Einkommensteuer um
\wert{kanal}{rv}{est}\,\Mrd, weil die Haushalte weniger Netto haben und weniger konsumieren.
Der Nachfrageimpuls sinkt um \wert{kanal}{rv}{nachfr}\,\Mrd. Unter dem Strich bleiben
\wert{p1}{rv}{saldo}\,\Mrd.

\newcommand{\ressortgraph}[3]{%
\begin{figure}[p]
\centering
\begin{tikzpicture}
\input{generiert/graph_#1.tex}
\end{tikzpicture}
\caption{#2}
\label{#3}
\end{figure}}

\ressortgraph{Finanzen}{Ressort Finanzen und Steuern: gemessene Wirkung auf die Kanäle,
erste Periode, je Schritt.}{fig:finanzen}
\ressortgraph{Wirtschaft}{Ressort Wirtschaft und Unternehmen. Die Investitionsreaktion wirkt
zusätzlich über das private Kapital auf das BIP der Folgeperioden
(Abbildung~\ref{fig:struktur}).}{fig:wirtschaft}
\ressortgraph{Soziales}{Ressort Arbeit und Soziales. Beiträge wirken nur auf die Einnahmen,
das Rentenniveau nur auf die Ausgaben.}{fig:soziales}
\ressortgraph{Umwelt}{Ressort Umwelt und Klima. Das Klimageld verschiebt das
\COz-Aufkommen vollständig zu den Transfers.}{fig:umwelt}

\paragraph{Was die Ressortgraphen zeigen.}
\begin{itemize}
\item \textbf{Jede Steuer bewegt mehr als ihren eigenen Kanal.} Eine höhere Umsatzsteuer
  senkt auch Einkommensteuer und Beiträge, weil sie das Vornetto nicht ändert, wohl aber das
  Konsumvolumen; eine höhere Beitragsrate senkt Einkommen- und Umsatzsteuer.
\item \textbf{Der Nachfrageimpuls ist fast immer die zweitgrößte Kante.} Er ist der Grund,
  warum keine Konsolidierung umsonst ist.
\item \textbf{Nur zwei Stellgrößen am Tisch bewegen das Arbeitsangebot} messbar:
  Grundfreibetrag und Eingangssatz, weil sie den Grenzsteuersatz der Mitte ändern. Der
  Spitzensatz trifft zu wenige Einkommen, und Beiträge gehen im Modell nicht in den
  Grenzsteuersatz ein (Abschnitt~\ref{sec:zerlegung}).
\end{itemize}

\section{Die Gesamtmatrix: Zielkonflikte auf einen Blick}
\label{sec:matrix}

Tabelle~\ref{tab:gesamt} fasst alle Pfade zusammen: Jede Zeile ist eine Stellgröße, jede
Spalte eine Kennzahl, jede Zelle die gemessene Gesamtwirkung. Die Färbung zeigt diesmal nicht
das Vorzeichen, sondern die \emph{Erwünschtheit}: \textcolor{pos}{\textbf{blau}} heißt, die
Kennzahl entwickelt sich in die erwünschte Richtung (höherer Saldo, niedrigerer Gini,
niedrigere Emissionen, höheres BIP, höheres Nettoeinkommen, niedrigere Schuldenquote),
\textcolor{neg}{\textbf{orange}} das Gegenteil. Die Sättigung folgt dem Betrag, relativ zum
größten Wert der Spalte.

\begin{table}[p]
\centering\scriptsize
\setlength{\tabcolsep}{3pt}
% Erzeugt von docs/modell/messung.js — nicht von Hand bearbeiten
\begin{tabular}{lrrrrrrrrrrr}
\toprule
Stellgröße & \rotatebox{70}{Saldo} & \rotatebox{70}{Gini} & \rotatebox{70}{Armutsrisiko} & \rotatebox{70}{Emissionen} & \rotatebox{70}{BIP} & \rotatebox{70}{Netto D1} & \rotatebox{70}{Netto D5} & \rotatebox{70}{Netto D10c} & \rotatebox{70}{Schuldenquote 2041} & \rotatebox{70}{BIP 2041} & \rotatebox{70}{CO\textsubscript{2} 2025–2040} \\
 & {\tiny Mrd. \euro{}/a} & {\tiny Gini-Pkt.} & {\tiny \%} & {\tiny Mt/a} & {\tiny Mrd. \euro{}} & {\tiny \euro{}/a} & {\tiny \euro{}/a} & {\tiny \euro{}/a} & {\tiny Pp.} & {\tiny Mrd. \euro{}} & {\tiny Mt} \\
\midrule
Grundfreibetrag & \cellcolor{neg!45}$-$6,9 & \cellcolor{pos!35}$-$0,13 & \cellcolor{neg!70}0,29 & \textcolor{grau}{0} & \cellcolor{pos!60}8,5 & \textcolor{grau}{0} & \cellcolor{pos!69}358 & \cellcolor{pos!24}637 & \cellcolor{neg!25}0,85 & \cellcolor{pos!70}28,2 & \textcolor{grau}{0} \\
Eingangssteuersatz & \cellcolor{pos!24}2,6 & \cellcolor{pos!19}$-$0,04 & \cellcolor{pos!29}$-$0,09 & \textcolor{grau}{0} & \cellcolor{neg!28}$-$2,8 & \textcolor{grau}{0} & \cellcolor{neg!28}$-$101 & \cellcolor{neg!18}$-$322 & \cellcolor{pos!20}$-$0,51 & \cellcolor{neg!28}$-$7,9 & \textcolor{grau}{0} \\
Spitzensteuersatz & \cellcolor{pos!13}0,3 & \cellcolor{pos!15}$-$0,02 & \textcolor{grau}{0} & \textcolor{grau}{0} & \cellcolor{neg!12}$-$0,1 & \textcolor{grau}{0} & \textcolor{grau}{0} & \cellcolor{neg!33}$-$1.156 & \cellcolor{pos!13}$-$0,09 & \cellcolor{neg!12}$-$0,2 & \textcolor{grau}{0} \\
Umsatzsteuer & \cellcolor{pos!49}7,7 & \cellcolor{neg!22}0,06 & \cellcolor{pos!19}$-$0,03 & \textcolor{grau}{0} & \cellcolor{neg!43}$-$5,5 & \cellcolor{neg!29}$-$87 & \cellcolor{neg!44}$-$204 & \cellcolor{neg!41}$-$1.555 & \cellcolor{pos!48}$-$2,44 & \cellcolor{neg!29}$-$8,2 & \textcolor{grau}{0} \\
Erbschaftsteuer & \cellcolor{pos!13}0,3 & \cellcolor{pos!13}$-$0,01 & \textcolor{grau}{0} & \textcolor{grau}{0} & \cellcolor{neg!13}$-$0,1 & \textcolor{grau}{0} & \cellcolor{neg!12}$-$2 & \cellcolor{neg!15}$-$177 & \cellcolor{pos!13}$-$0,08 & \cellcolor{neg!12}$-$0,2 & \textcolor{grau}{0} \\
Körperschaftsteuer & \cellcolor{pos!22}2,0 & \cellcolor{pos!18}$-$0,04 & \cellcolor{pos!14}$-$0,01 & \textcolor{grau}{0} & \cellcolor{neg!18}$-$1,0 & \cellcolor{neg!14}$-$8 & \cellcolor{neg!16}$-$27 & \cellcolor{neg!42}$-$1.616 & \cellcolor{pos!17}$-$0,33 & \cellcolor{neg!29}$-$8,1 & \textcolor{grau}{0} \\
Gewerbesteuer & \cellcolor{pos!31}3,9 & \cellcolor{pos!25}$-$0,07 & \cellcolor{pos!16}$-$0,02 & \textcolor{grau}{0} & \cellcolor{neg!23}$-$2,0 & \cellcolor{neg!15}$-$16 & \cellcolor{neg!20}$-$52 & \cellcolor{neg!70}$-$3.128 & \cellcolor{pos!26}$-$0,93 & \cellcolor{neg!32}$-$9,6 & \textcolor{grau}{0} \\
RV-Beitrag & \cellcolor{pos!68}11,7 & \cellcolor{neg!24}0,06 & \cellcolor{pos!33}$-$0,10 & \textcolor{grau}{0} & \cellcolor{neg!70}$-$10,3 & \cellcolor{neg!36}$-$129 & \cellcolor{neg!70}$-$366 & \cellcolor{neg!30}$-$987 & \cellcolor{pos!68}$-$3,76 & \cellcolor{neg!44}$-$15,6 & \textcolor{grau}{0} \\
KV-Beitrag & \cellcolor{pos!70}12,1 & \cellcolor{neg!34}0,12 & \cellcolor{pos!33}$-$0,10 & \textcolor{grau}{0} & \cellcolor{neg!66}$-$9,6 & \cellcolor{neg!36}$-$129 & \cellcolor{neg!70}$-$366 & \cellcolor{neg!25}$-$679 & \cellcolor{pos!70}$-$3,90 & \cellcolor{neg!42}$-$14,5 & \textcolor{grau}{0} \\
Rentenniveau & \cellcolor{neg!40}$-$5,8 & \cellcolor{pos!51}$-$0,22 & \cellcolor{pos!13}$-$0,01 & \textcolor{grau}{0} & \cellcolor{pos!39}4,8 & \cellcolor{pos!35}122 & \cellcolor{pos!50}238 & \cellcolor{pos!14}117 & \cellcolor{neg!42}2,01 & \cellcolor{pos!31}9,0 & \textcolor{grau}{0} \\
Bürgergeld & \cellcolor{neg!14}$-$0,4 & \cellcolor{pos!20}$-$0,04 & \cellcolor{pos!27}$-$0,08 & \textcolor{grau}{0} & \cellcolor{pos!14}0,4 & \cellcolor{pos!26}72 & \textcolor{grau}{0} & \textcolor{grau}{0} & \cellcolor{neg!14}0,11 & \cellcolor{pos!13}0,6 & \textcolor{grau}{0} \\
Kindergeld & \cellcolor{neg!20}$-$1,7 & \cellcolor{pos!24}$-$0,07 & \cellcolor{pos!16}$-$0,02 & \textcolor{grau}{0} & \cellcolor{pos!19}1,3 & \cellcolor{pos!20}45 & \cellcolor{pos!21}59 & \cellcolor{pos!12}11 & \cellcolor{neg!20}0,52 & \cellcolor{pos!16}2,0 & \textcolor{grau}{0} \\
CO\textsubscript{2}-Preis & \cellcolor{pos!22}2,0 & \cellcolor{neg!16}0,02 & \cellcolor{pos!14}$-$0,01 & \cellcolor{pos!70}$-$9,8 & \cellcolor{neg!20}$-$1,4 & \cellcolor{neg!17}$-$24 & \cellcolor{neg!21}$-$55 & \cellcolor{neg!18}$-$303 & \cellcolor{pos!19}$-$0,46 & \cellcolor{neg!14}$-$0,8 & \cellcolor{pos!70}$-$103 \\
Klimageld & \cellcolor{neg!58}$-$9,7 & \cellcolor{pos!70}$-$0,32 & \cellcolor{pos!43}$-$0,15 & \textcolor{grau}{0} & \cellcolor{pos!56}7,8 & \cellcolor{pos!70}307 & \cellcolor{pos!61}307 & \cellcolor{pos!18}307 & \cellcolor{neg!45}2,19 & \cellcolor{pos!21}4,3 & \textcolor{grau}{0} \\
\bottomrule
\end{tabular}

\caption{Gesamtmatrix der Wirkungen. Die ersten acht Spalten gelten für die erste Periode
(Jahreswerte 2025--2028), die letzten drei für das Ende des Pfads nach fünf Perioden mit
unveränderter Politik. \textcolor{pos}{Blau}: erwünschte Richtung, \textcolor{neg}{orange}:
unerwünschte. Grau: keine messbare Wirkung. Schritte wie in Tabelle~\ref{tab:schritte}.}
\label{tab:gesamt}
\end{table}

\paragraph{So liest man die Matrix.} Lesen Sie eine \emph{Zeile}: Sie zeigt den Preis einer
Stellgröße. Keine Zeile ist durchgehend blau, das ist die Eigenschaft (E5) aus
Abschnitt~\ref{sec:eigenschaften} und wird von \datei{tests/dominanz.test.js} erzwungen. Beispiel
\emph{Grundfreibetrag}: Er entlastet D5 und D10c, senkt den Gini und hebt das BIP, kostet aber
\wert{p1}{freibetrag}{saldo}\,\Mrd{} Saldo und erhöht die Schuldenquote 2041 um
\wert{pfad}{freibetrag}{schuld}\,\Pp. Dass er die Armutsrisikoquote \emph{erhöht}
(\wert{p1}{freibetrag}{armut}\,\Pp), ist kein Fehler: D1 liegt schon unter dem Freibetrag und
gewinnt nichts, die Mitte gewinnt, das mittlere Einkommen steigt, und mit ihm die relative
Armutsgrenze.

Lesen Sie eine \emph{Spalte}: Sie zeigt, welche Hebel eine Kennzahl überhaupt bewegen. Die
Emissionen bewegt allein der \COz-Preis; das BIP 2041 bewegen am stärksten Freibetrag,
Beiträge und Rentenniveau, über Nachfrage, Arbeitsangebot und Schulden.

\section{Verteilung über die Dezile}

Tabelle~\ref{tab:dezile-wirkung} zeigt die Änderung des Nettoeinkommens jedes Haushaltstyps
in \euro{} im Jahr. Sie macht die Inzidenz sichtbar, die in Gini und Palma nur verdichtet
erscheint.

\begin{table}[!htb]
\centering\scriptsize
\setlength{\tabcolsep}{2.5pt}
% Erzeugt von docs/modell/messung.js — nicht von Hand bearbeiten
\begin{tabular}{lrrrrrrrrrrrr}
\toprule
 & D1 & D2 & D3 & D4 & D5 & D6 & D7 & D8 & D9 & D10a & D10b & D10c \\
\midrule
Grundfreibetrag & 0 & 0 & $+$642 & $+$341 & $+$358 & $+$381 & $+$409 & $+$447 & $+$515 & $+$613 & $+$626 & $+$637 \\
Eingangssteuersatz & 0 & 0 & $-$36 & $-$71 & $-$101 & $-$137 & $-$176 & $-$223 & $-$284 & $-$327 & $-$317 & $-$322 \\
Spitzensteuersatz & 0 & 0 & 0 & 0 & 0 & 0 & 0 & 0 & 0 & 0 & 0 & $-$1.156 \\
Umsatzsteuer & $-$87 & $-$130 & $-$161 & $-$182 & $-$204 & $-$229 & $-$260 & $-$298 & $-$361 & $-$397 & $-$618 & $-$1.555 \\
Erbschaftsteuer & 0 & 0 & 0 & $-$1 & $-$2 & $-$3 & $-$5 & $-$9 & $-$16 & $-$21 & $-$38 & $-$177 \\
Körperschaftsteuer & $-$8 & $-$12 & $-$17 & $-$21 & $-$27 & $-$32 & $-$41 & $-$54 & $-$78 & $-$108 & $-$276 & $-$1.616 \\
Gewerbesteuer & $-$16 & $-$24 & $-$33 & $-$40 & $-$52 & $-$62 & $-$79 & $-$104 & $-$152 & $-$209 & $-$534 & $-$3.128 \\
RV-Beitrag & $-$129 & $-$194 & $-$248 & $-$304 & $-$366 & $-$440 & $-$527 & $-$649 & $-$842 & $-$974 & $-$979 & $-$987 \\
KV-Beitrag & $-$129 & $-$194 & $-$248 & $-$304 & $-$366 & $-$440 & $-$527 & $-$649 & $-$665 & $-$670 & $-$674 & $-$679 \\
Rentenniveau & $+$122 & $+$202 & $+$236 & $+$243 & $+$238 & $+$218 & $+$182 & $+$152 & $+$133 & $+$109 & $+$115 & $+$117 \\
Bürgergeld & $+$72 & $+$30 & $+$10 & $+$2 & 0 & 0 & 0 & 0 & 0 & 0 & 0 & 0 \\
Kindergeld & $+$45 & $+$61 & $+$67 & $+$64 & $+$59 & $+$53 & $+$47 & $+$42 & $+$36 & $+$28 & $+$20 & $+$11 \\
CO\textsubscript{2}-Preis & $-$24 & $-$35 & $-$42 & $-$49 & $-$55 & $-$62 & $-$70 & $-$78 & $-$90 & $-$97 & $-$143 & $-$303 \\
Klimageld & $+$307 & $+$307 & $+$307 & $+$307 & $+$307 & $+$307 & $+$307 & $+$307 & $+$307 & $+$307 & $+$307 & $+$307 \\
\bottomrule
\end{tabular}

\caption{Änderung des Nettoeinkommens je Haushalt und Jahr in \euro, erste Periode, je
Schritt. D10a bis D10c teilen das oberste Dezil (90.--95., 95.--99., oberstes Prozent).}
\label{tab:dezile-wirkung}
\end{table}

\paragraph{So liest man die Tabelle.} Vergleichen Sie Beträge innerhalb einer Zeile mit den
Einkommen aus Tabelle~\ref{tab:dezile}: Die \COz-Abgabe kostet D1 \wert{p1}{co2}{d1}\,\euro{}
und D10c \wert{p1}{co2}{d10c}\,\euro{}; gemessen am Einkommen (14.000 gegen 700.000\,\euro)
trifft sie D1 aber viermal so stark. Das Klimageld zahlt allen gleich viel und kehrt diese
Last deshalb um.

\section{Wirkung über den ganzen Reglerbereich}
\label{sec:kurven}

Die Gewichte oben gelten für einen Schritt am Status quo. Abbildung~\ref{fig:kurven} zeigt
für fünf Stellgrößen, wie Saldo und Verteilung (beim \COz-Preis: Emissionen) über den
ganzen Bereich des Reglers verlaufen; alle anderen Regler bleiben im Status quo. Die
gestrichelte Linie markiert den Status quo.

\newcommand{\kurve}[6]{% Stellgröße, Spalte, x-Beschriftung, y-Beschriftung, Titel, Farbe
  \nextgroupplot[title={#5}, xlabel={#3}, ylabel={#4},
    extra x ticks={\wert{sq}{#1}{x}}, extra x tick labels={},
    extra x tick style={grid=major, major grid style={dashed, grau}}]
  \addplot[color=#6, very thick, mark=none] table[x=x, y=#2] {generiert/verlauf_#1.dat};}

\begin{figure}[p]
\centering
\begin{tikzpicture}
\begin{groupplot}[group style={group size=2 by 5, horizontal sep=18mm, vertical sep=21mm},
  width=0.44\textwidth, height=0.15\textheight,
  tick label style={font=\scriptsize}, label style={font=\scriptsize},
  title style={font=\small\bfseries}, grid=none,
  /pgf/number format/use comma, /pgf/number format/1000 sep={.},
  scaled y ticks=false]
\kurve{spitze}{saldo}{Spitzensteuersatz [\%]}{Saldo [\Mrd]}{Spitzensteuersatz: Saldo}{pos}
\kurve{spitze}{gini}{Spitzensteuersatz [\%]}{Gini [Pkt.]}{Spitzensteuersatz: Gini}{neg}
\kurve{mwst}{saldo}{Umsatzsteuer [\%]}{Saldo [\Mrd]}{Umsatzsteuer: Saldo}{pos}
\kurve{mwst}{gini}{Umsatzsteuer [\%]}{Gini [Pkt.]}{Umsatzsteuer: Gini}{neg}
\kurve{kst}{saldo}{Körperschaftsteuer [\%]}{Saldo [\Mrd]}{Körperschaftsteuer: Saldo}{pos}
\kurve{kst}{bip}{Körperschaftsteuer [\%]}{BIP [\Mrd]}{Körperschaftsteuer: BIP der Periode}{neg}
\kurve{rentenniveau}{saldo}{Rentenniveau [\%]}{Saldo [\Mrd]}{Rentenniveau: Saldo}{pos}
\kurve{rentenniveau}{gini}{Rentenniveau [\%]}{Gini [Pkt.]}{Rentenniveau: Gini}{neg}
\kurve{co2}{saldo}{\COz-Preis [\euro/t]}{Saldo [\Mrd]}{\COz-Preis: Saldo}{pos}
\kurve{co2}{emiss}{\COz-Preis [\euro/t]}{Emissionen [Mt/a]}{\COz-Preis: Emissionen}{neg}
\end{groupplot}
\end{tikzpicture}
\caption{Saldo und eine zweite Kennzahl der ersten Periode über den ganzen Bereich des
Reglers, alle übrigen Regler im Status quo. Gestrichelt: Wert des Status quo. 21 Stützstellen
je Kurve.}
\label{fig:kurven}
\end{figure}

\paragraph{So liest man die Kurven.} Eine Gerade heißt: Jeder Schritt wirkt gleich, die
Gewichte oben gelten im ganzen Bereich. Ein Knick heißt: Dort greift eine Schwelle oder
Klammer der Engine (Abschnitt~\ref{sec:eigenschaften}). Eine flacher werdende Kurve heißt:
Die Verhaltensreaktion frisst einen wachsenden Teil der mechanischen Wirkung. Beispiele:
\begin{itemize}
\item \textbf{Spitzensteuersatz:} Unterhalb von 42\,\% wird die Proportionalzone auf den
  Spitzensatz gedeckelt, der Satz trifft nun die ganze obere Mitte, und der Saldo fällt
  steil. Oberhalb wirkt er nur im Pareto-Rand des obersten Prozents, und die Kurve ist flach
  und gerade. Die Ausweich- und Wegzugsreaktion des obersten Prozents setzt im ganzen
  Bereich \emph{nicht} ein (Abschnitt~\ref{sec:grenzen}).
\item \textbf{Umsatzsteuer:} leicht konkav, weil die Konsumreaktion mit dem Satz wächst und
  die Nachfrage mitzieht; eine Obergrenze erreicht sie im Bereich nicht.
\item \textbf{Körperschaftsteuer:} Das Aufkommen ist Satz mal schrumpfende Basis, also eine
  nach unten geöffnete Parabel; ihr Scheitel liegt weit außerhalb des Reglerbereichs, die
  Klammern der Investitionsreaktion werden nicht erreicht.
\item \textbf{\COz-Preis:} Die Emissionen fallen linear zwischen etwa 22\,\euro/t (darunter
  greift die obere Klammer) und etwa 255\,\euro/t (darüber greift die untere Klammer bei
  40\,\% der bepreisten Emissionen). Der Saldo zeigt eine \emph{Laffer-Kurve}: Das Aufkommen
  $p\cdot E(p)$ erreicht sein Maximum bei etwa 190\,\euro/t, weil die Bemessungsgrundlage
  schneller schrumpft, als der Preis steigt; jenseits der unteren Klammer steigt es wieder.
\end{itemize}

\section{Wirkung über den Pfad}

Tabelle~\ref{tab:pfad} zeigt, was nach fünf Perioden mit unveränderter Politik bleibt. Hier
wirken die Kanäle, die in der ersten Periode unsichtbar sind: das private Kapital, das
Arbeitsangebot im Niveau und vor allem der Zinseszins auf die Schulden.

\begin{table}[!htb]
\centering\small
% Erzeugt von docs/modell/messung.js — nicht von Hand bearbeiten
\begin{tabular}{lrrrrr}
\toprule
Stellgröße (Schritt) & \rotatebox{75}{Schuldenquote 2041 [Pp.]} & \rotatebox{75}{BIP 2041 [Mrd. \euro{}]} & \rotatebox{75}{CO\textsubscript{2} 2025–2040 [Mt]} & \rotatebox{75}{Saldo 2041–44 [Mrd. \euro{}/a]} & \rotatebox{75}{Gini 2041–44 [Gini-Pkt.]} \\
\midrule
Grundfreibetrag (1.000 \euro{}) & $+$0,85 & $+$28,2 & 0 & $-$3,6 & $-$0,124 \\
Eingangssteuersatz (1 Pp.) & $-$0,51 & $-$7,9 & 0 & $+$2,7 & $-$0,037 \\
Spitzensteuersatz (1 Pp.) & $-$0,09 & $-$0,2 & 0 & $+$0,6 & $-$0,018 \\
Umsatzsteuer (1 Pp.) & $-$2,44 & $-$8,2 & 0 & $+$14,7 & $+$0,057 \\
Erbschaftsteuer (1 Pp.) & $-$0,08 & $-$0,2 & 0 & $+$0,5 & $-$0,007 \\
Körperschaftsteuer (1 Pp.) & $-$0,33 & $-$8,1 & 0 & $+$1,4 & $-$0,038 \\
Gewerbesteuer (1 Pp.) & $-$0,93 & $-$9,6 & 0 & $+$5,0 & $-$0,072 \\
RV-Beitrag (1 Pp.) & $-$3,76 & $-$15,6 & 0 & $+$22,8 & $+$0,066 \\
KV-Beitrag (1 Pp.) & $-$3,90 & $-$14,5 & 0 & $+$23,6 & $+$0,122 \\
Rentenniveau (1 Pp.) & $+$2,01 & $+$9,0 & 0 & $-$13,3 & $-$0,275 \\
Bürgergeld (10 \euro{}/Monat) & $+$0,11 & $+$0,6 & 0 & $-$0,7 & $-$0,043 \\
Kindergeld (10 \euro{}/Monat) & $+$0,52 & $+$2,0 & 0 & $-$3,1 & $-$0,067 \\
CO\textsubscript{2}-Preis (10 \euro{}/t) & $-$0,46 & $-$0,8 & $-$103 & $+$1,7 & $+$0,007 \\
Klimageld (ein) & $+$2,19 & $+$4,3 & 0 & $-$8,2 & $-$0,119 \\
\bottomrule
\end{tabular}

\caption{Wirkung nach fünf Perioden (Ende des Pfads) gegenüber dem Status-quo-Pfad, je
Schritt, Politik über alle Perioden unverändert. Status quo: Schuldenquote \SQschuldEnde\,\%,
BIP \SQbipEnde\,\Mrd, \COz{} kumuliert \SQcoEnde\,Mt.}
\label{tab:pfad}
\end{table}

\paragraph{So liest man die Tabelle.} Teilen Sie die Änderung der Schuldenquote durch die
Saldowirkung der ersten Periode. Wirkt eine Stellgröße nur über den Saldo, liegt dieses
Verhältnis bei rund $-0{,}3$\,\Pp{} je Mrd.\,\euro{} (etwa Rentenbeitrag, Umsatzsteuer,
Kindergeld): Ein dauerhafter Saldo geht über 16 Jahre fast linear in die Quote ein. Weicht
das Verhältnis ab, ist ein weiterer Kanal im Spiel, der das BIP im Nenner bewegt:
\begin{itemize}
\item Der \textbf{Grundfreibetrag} kostet \wert{p1}{freibetrag}{saldo}\,\Mrd{} im Jahr, erhöht
  die Quote 2041 aber nur um \wert{pfad}{freibetrag}{schuld}\,\Pp: Das zusätzliche
  Arbeitsangebot hebt das BIP 2041 um \wert{pfad}{freibetrag}{bip}\,\Mrd.
\item Die \textbf{Körperschaftsteuer} bringt \wert{p1}{kst}{saldo}\,\Mrd, senkt die Quote aber
  nur um \wert{pfad}{kst}{schuld}\,\Pp: Das private Kapital schrumpft, das BIP 2041 sinkt um
  \wert{pfad}{kst}{bip}\,\Mrd.
\end{itemize}

\section{Die Messwerte im Einzelnen}

Zum Nachschlagen die vollständigen Messwerte für die Stellgrößen am Tisch.

\begin{table}[!htb]
\centering\scriptsize
\setlength{\tabcolsep}{3pt}
% Erzeugt von docs/modell/messung.js — nicht von Hand bearbeiten
\begin{tabular}{lrrrrrrrrrrr}
\toprule
Stellgröße (Schritt) & \rotatebox{75}{Einkommensteuer [Mrd. \euro{}]} & \rotatebox{75}{Umsatzsteuer [Mrd. \euro{}]} & \rotatebox{75}{KSt + GewSt [Mrd. \euro{}]} & \rotatebox{75}{Erb-, Boden-, Verm.-St. [Mrd. \euro{}]} & \rotatebox{75}{SV-Beiträge [Mrd. \euro{}]} & \rotatebox{75}{CO\textsubscript{2}-Einnahmen netto [Mrd. \euro{}]} & \rotatebox{75}{Transfers [Mrd. \euro{}]} & \rotatebox{75}{Rentenzahlungen [Mrd. \euro{}]} & \rotatebox{75}{Arbeitsangebot [\%]} & \rotatebox{75}{Investitionsfaktor [\%]} & \rotatebox{75}{Nachfrageimpuls [Mrd. \euro{}]} \\
\midrule
Grundfreibetrag (1.000 \euro{}) & $-$12,6 & $+$3,4 & $+$0,2 & $+$0,1 & $+$1,4 & 0,0 & 0,0 & 0,0 & $+$0,35 & 0,00 & $+$8,5 \\
Eingangssteuersatz (1 Pp.) & $+$4,5 & $-$1,1 & $-$0,1 & 0,0 & $-$0,5 & 0,0 & 0,0 & 0,0 & $-$0,08 & 0,00 & $-$2,8 \\
Spitzensteuersatz (1 Pp.) & $+$0,3 & 0,0 & 0,0 & 0,0 & 0,0 & 0,0 & 0,0 & 0,0 & 0,00 & 0,00 & $-$0,1 \\
Umsatzsteuer (1 Pp.) & $-$0,5 & $+$9,5 & $-$0,1 & 0,0 & $-$0,9 & 0,0 & 0,0 & 0,0 & 0,00 & 0,00 & $-$5,5 \\
Erbschaftsteuer (1 Pp.) & 0,0 & 0,0 & 0,0 & $+$0,3 & 0,0 & 0,0 & 0,0 & 0,0 & 0,00 & 0,00 & $-$0,1 \\
Körperschaftsteuer (1 Pp.) & $-$0,1 & $-$0,1 & $+$2,5 & 0,0 & $-$0,2 & 0,0 & 0,0 & 0,0 & 0,00 & $-$0,40 & $-$1,0 \\
Gewerbesteuer (1 Pp.) & $-$0,2 & $-$0,1 & $+$4,8 & 0,0 & $-$0,3 & 0,0 & 0,0 & 0,0 & 0,00 & $-$0,40 & $-$2,0 \\
RV-Beitrag (1 Pp.) & $-$0,9 & $-$2,3 & $-$0,3 & $-$0,1 & $+$15,7 & 0,0 & 0,0 & 0,0 & 0,00 & 0,00 & $-$10,3 \\
KV-Beitrag (1 Pp.) & $-$0,8 & $-$2,1 & $-$0,3 & $-$0,1 & $+$15,9 & 0,0 & 0,0 & 0,0 & 0,00 & 0,00 & $-$9,6 \\
Rentenniveau (1 Pp.) & $+$0,4 & $+$0,3 & $+$0,1 & 0,0 & $+$0,8 & 0,0 & 0,0 & $+$7,5 & 0,00 & 0,00 & $+$4,8 \\
Bürgergeld (10 \euro{}/Monat) & 0,0 & 0,0 & 0,0 & 0,0 & $+$0,1 & 0,0 & $+$0,5 & 0,0 & 0,00 & 0,00 & $+$0,4 \\
Kindergeld (10 \euro{}/Monat) & $+$0,1 & $+$0,1 & 0,0 & 0,0 & $+$0,2 & 0,0 & $+$2,0 & 0,0 & 0,00 & 0,00 & $+$1,3 \\
CO\textsubscript{2}-Preis (10 \euro{}/t) & $-$0,1 & $-$0,1 & 0,0 & 0,0 & $-$0,2 & $+$2,6 & 0,0 & 0,0 & 0,00 & 0,00 & $-$1,4 \\
Klimageld (ein) & $+$0,6 & $+$0,5 & $+$0,2 & $+$0,1 & $+$1,3 & $-$12,6 & $+$12,6 & 0,0 & 0,00 & 0,00 & $+$7,8 \\
\bottomrule
\end{tabular}

\caption{Wirkung auf die Kanäle (erste Periode, je Schritt). Das sind die Kantengewichte der
Ressortgraphen, einschließlich der dort weggelassenen kleinen Werte.}
\label{tab:kanaele}
\end{table}

\begin{table}[!htb]
\centering\scriptsize
\setlength{\tabcolsep}{3pt}
% Erzeugt von docs/modell/messung.js — nicht von Hand bearbeiten
\begin{tabular}{lrrrrrrrr}
\toprule
Stellgröße (Schritt) & \rotatebox{75}{Saldo [Mrd. \euro{}/a]} & \rotatebox{75}{Gini [Gini-Pkt.]} & \rotatebox{75}{Armutsrisiko [\%]} & \rotatebox{75}{Emissionen [Mt/a]} & \rotatebox{75}{BIP [Mrd. \euro{}]} & \rotatebox{75}{Netto D1 [\euro{}/a]} & \rotatebox{75}{Netto D5 [\euro{}/a]} & \rotatebox{75}{Netto D10c [\euro{}/a]} \\
\midrule
Grundfreibetrag (1.000 \euro{}) & $-$6,9 & $-$0,125 & $+$0,29 & 0,0 & $+$8,5 & 0 & $+$358 & $+$637 \\
Eingangssteuersatz (1 Pp.) & $+$2,6 & $-$0,036 & $-$0,09 & 0,0 & $-$2,8 & 0 & $-$101 & $-$322 \\
Spitzensteuersatz (1 Pp.) & $+$0,3 & $-$0,017 & 0,00 & 0,0 & $-$0,1 & 0 & 0 & $-$1.156 \\
Umsatzsteuer (1 Pp.) & $+$7,7 & $+$0,057 & $-$0,03 & 0,0 & $-$5,5 & $-$87 & $-$204 & $-$1.555 \\
Erbschaftsteuer (1 Pp.) & $+$0,3 & $-$0,007 & 0,00 & 0,0 & $-$0,1 & 0 & $-$2 & $-$177 \\
Körperschaftsteuer (1 Pp.) & $+$2,0 & $-$0,035 & $-$0,01 & 0,0 & $-$1,0 & $-$8 & $-$27 & $-$1.616 \\
Gewerbesteuer (1 Pp.) & $+$3,9 & $-$0,068 & $-$0,02 & 0,0 & $-$2,0 & $-$16 & $-$52 & $-$3.128 \\
RV-Beitrag (1 Pp.) & $+$11,7 & $+$0,064 & $-$0,10 & 0,0 & $-$10,3 & $-$129 & $-$366 & $-$987 \\
KV-Beitrag (1 Pp.) & $+$12,1 & $+$0,119 & $-$0,10 & 0,0 & $-$9,6 & $-$129 & $-$366 & $-$679 \\
Rentenniveau (1 Pp.) & $-$5,8 & $-$0,216 & $-$0,01 & 0,0 & $+$4,8 & $+$122 & $+$238 & $+$117 \\
Bürgergeld (10 \euro{}/Monat) & $-$0,4 & $-$0,042 & $-$0,08 & 0,0 & $+$0,4 & $+$72 & 0 & 0 \\
Kindergeld (10 \euro{}/Monat) & $-$1,7 & $-$0,066 & $-$0,02 & 0,0 & $+$1,3 & $+$45 & $+$59 & $+$11 \\
CO\textsubscript{2}-Preis (10 \euro{}/t) & $+$2,0 & $+$0,019 & $-$0,01 & $-$9,8 & $-$1,4 & $-$24 & $-$55 & $-$303 \\
Klimageld (ein) & $-$9,7 & $-$0,317 & $-$0,15 & 0,0 & $+$7,8 & $+$307 & $+$307 & $+$307 \\
\bottomrule
\end{tabular}

\caption{Wirkung auf die Kennzahlen (erste Periode, je Schritt).}
\label{tab:kennzahlen}
\end{table}

\begin{table}[!htb]
\centering\scriptsize
\setlength{\tabcolsep}{3pt}
% Erzeugt von docs/modell/messung.js — nicht von Hand bearbeiten
\begin{tabular}{lrrrrrrrr}
\toprule
Stellgröße (Schritt) & \rotatebox{75}{Saldo [Mrd. \euro{}/a]} & \rotatebox{75}{Gini [Gini-Pkt.]} & \rotatebox{75}{Armutsrisiko [\%]} & \rotatebox{75}{Emissionen [Mt/a]} & \rotatebox{75}{BIP [Mrd. \euro{}]} & \rotatebox{75}{Netto D1 [\euro{}/a]} & \rotatebox{75}{Netto D5 [\euro{}/a]} & \rotatebox{75}{Netto D10c [\euro{}/a]} \\
\midrule
Grenze Spitzensatz ($-$10.000 \euro{}) & 0,0 & $-$0,001 & 0,00 & 0,0 & 0,0 & 0 & 0 & $-$92 \\
USt ermäßigt (1 Pp.) & $+$4,4 & $+$0,030 & $-$0,02 & 0,0 & $-$2,9 & $-$46 & $-$108 & $-$824 \\
Abgeltungsteuer (1 Pp.) & $+$2,3 & $-$0,072 & $-$0,01 & 0,0 & $-$0,8 & $-$1 & $-$11 & $-$2.984 \\
Beitragsbemessungsgr. (10.000 \euro{}) & $+$6,0 & $-$0,274 & 0,00 & 0,0 & $-$3,6 & 0 & 0 & $-$1.704 \\
Vermögensteuer (0,1 Pp.) & $+$2,8 & $-$0,075 & $-$0,02 & 0,0 & $-$1,3 & 0 & $-$19 & $-$1.868 \\
Bodenwertsteuer (0,1 Pp.) & $+$4,2 & $-$0,109 & $-$0,02 & 0,0 & $-$2,0 & 0 & $-$27 & $-$2.726 \\
Zucman-Mindeststeuer (1 Pp.) & $+$3,4 & $-$0,142 & 0,00 & 0,0 & $-$0,7 & 0 & 0 & $-$9.476 \\
\bottomrule
\end{tabular}

\caption{Wirkung der Stellgrößen der klassischen Fläche auf die Kennzahlen (erste Periode).}
\label{tab:weitere}
\end{table}
